\section{Methodology}

This section describes the datasets, graph construction pipeline, feature extraction, anomaly scoring methodology, and experimental setup used to evaluate our approach. Our central research question is whether combining network flow logs and authentication logs into a unified graph representation yields better lateral movement detection than either source alone, and whether this graph-based approach outperforms classical anomaly detection methods operating on tabular features.

\subsection{Datasets}

We evaluate our approach on two public cybersecurity datasets that provide complementary views of the lateral movement detection problem.

\textbf{LANL Cybersecurity Dataset (LANL-2015).}
The primary dataset is the Los Alamos National Laboratory's comprehensive cybersecurity dataset from 2015~\cite{lanl2015}, which captures 58 days of continuous network activity across an enterprise network. It contains approximately 1.6 billion events spanning 12,425 unique users and 17,684 unique computers, organized into several gzipped CSV files: \texttt{auth.txt.gz} records authentication events (timestamp, source and destination users and computers, authentication type, logon type, orientation, and success/failure status); \texttt{flows.txt.gz} records network flow events (timestamp, duration, source and destination computers and ports, protocol, packet count, and byte count); and \texttt{redteam.txt.gz} provides ground-truth labels for 749 red team attack events spanning 308 unique source-destination host pairs. The red team events include lateral movement via SMB (port 445), SSH (port 22), RDP (port 3389), credential reuse, and network reconnaissance, with 93.6\% of attacks originating from a single compromised host serving as a jump box.

To scope the data to computationally tractable and analytically relevant periods, we apply a time-windowing strategy. For each of the 749 red team events, we define a window spanning $\pm$3,600 seconds around the event timestamp. Overlapping windows are merged, producing 25 non-overlapping intervals that capture approximately 104 million authentication events and 31.6 million flow events. Because the red team events span a large portion of the dataset's time range, these merged windows cover most of the activity while excluding long quiet periods.

\textbf{DAPT2020 Dataset.}
The secondary dataset is the Detecting and Preventing Cyber Attacks (DAPT2020) dataset~\cite{dapt2020}, which provides pre-extracted CICFlowMeter features from network traffic captures of simulated APT attack scenarios. It contains 86,690 flow records with 77 numerical features including packet length statistics, inter-arrival times, flag counts, and segment size averages, along with activity and stage labels. Lateral movement samples are identified by filtering the \texttt{Stage} column for entries containing ``lateral movement.'' This dataset serves as a benchmark for comparing our graph-based approach against classical unsupervised anomaly detection methods that operate on tabular features rather than graph structure.

\subsection{Graph Construction}

We model the network as a directed multigraph using the igraph library, where nodes represent entities (computers and users) and edges represent interactions (authentication events and network flows). The graph is constructed incrementally via a streaming approach: raw gzipped CSV files are read line-by-line, each event is validated against the time windows, and qualifying events are fed directly into the graph without materializing intermediate DataFrames in memory. This design enables the pipeline to process over 135 million events within approximately 4 GB of RAM.

For authentication events, the system creates a computer-to-computer edge from the source computer to the destination computer, and optionally a user-to-user edge from the source user to the destination user. Each edge stores attributes including authentication type, logon type, orientation, success status, and timestamp. When multiple events share the same source-destination pair, the edge weight is incremented and the first and last timestamps are preserved, producing a deduplicated graph with frequency information. For network flow events, a computer-to-computer edge is created for each source-destination pair, storing protocol, source and destination ports, packet count, byte count, duration, and timestamp. Node types are inferred heuristically: any node name ending with \texttt{\$} or containing \texttt{\$} before an \texttt{@} symbol is classified as a machine account, while others are classified as users.

Three graph variants are constructed for each experiment run: a \emph{flow-only} graph containing only network flow edges (approximately 11,586 nodes and 131,562 edges), an \emph{auth-only} graph containing only authentication edges (approximately 90,783 nodes and 409,100 edges), and a \emph{combined} graph incorporating both data sources (approximately 91,589 nodes and 512,529 edges). Each graph is freed from memory after scoring and evaluation, keeping peak memory bounded.

\subsection{Feature Extraction}

Once the graph is constructed, we extract structural, behavioral, and temporal features at the node, edge, and graph levels. All features are computed independently for each of the three graph variants.

\textbf{Node features.}
For each node, we compute in-degree, out-degree, total degree, and fan-out ratio (out-degree divided by total degree). Betweenness centrality is computed using exact igraph computation for graphs with fewer than 5,000 nodes; for larger graphs, we use an approximation with a path length cutoff of 3 hops, which captures local bridging behavior without the $O(n^3)$ cost of exact computation. Temporal features include mean and standard deviation of inter-arrival times between outgoing edges, a burst score measuring the maximum fraction of outgoing edges that fall within any 10\% window of the node's active time span, and total active duration.

\textbf{Edge features.}
For each edge, we compute edge rarity (the reciprocal of edge weight, measuring how unusual the connection frequency is), source out-degree and destination in-degree (contextualizing the edge within the local graph structure), and binary indicators for authentication-specific signals: whether the authentication uses NTLM (a protocol often associated with lateral movement), whether the logon type is Network (indicating remote access), and whether the authentication was successful. For flow edges, we additionally compute an unusual destination port indicator flagging ports commonly used in lateral movement (SSH-22, Telnet-23, SMB-445, RDP-3389, WinRM-5985/5986) or ephemeral ports above 49,152; protocol rarity measured as one minus the ratio of edges using that protocol to total flow edges; byte-per-packet ratio as a percentile rank across all edges; and duration z-score across the edge population. Self-loops and user-to-user edges are flagged and excluded from downstream anomaly scoring.

\textbf{Graph features.}
At the graph level, we compute density, average clustering coefficient, number of connected components, node count, and edge count. These features provide context for interpreting the structural complexity of each graph variant.

\subsection{Anomaly Scoring}

\subsubsection{Edge Scoring}
Edges are scored using a weighted combination of discriminative features, with the scoring formula differing by edge type. Authentication edges are scored as:
\begin{equation}
s_{\text{auth}} = 0.4 \cdot \text{is\_ntlm} + 0.3 \cdot \text{is\_network\_logon} + 0.3 \cdot \text{rarity\_rank}
\end{equation}
where \texttt{rarity\_rank} is the percentile rank of the edge's rarity value across all edges. This formula assigns higher scores to edges exhibiting known lateral movement indicators: NTLM authentication (uncommon in normal traffic but frequently used in pass-the-hash attacks), network logon type (indicating remote access rather than interactive sessions), and rare connections (low-frequency source-destination pairs).

Flow edges use a separate three-feature formula:
\begin{equation}
s_{\text{flow}} = 0.4 \cdot \text{rarity\_rank} + 0.3 \cdot \text{is\_unusual\_dst\_port} + 0.3 \cdot \text{protocol\_rarity\_rank}
\end{equation}
This gives flow edges meaningful anomaly signals beyond simple rarity, incorporating the suspiciousness of the destination port and protocol.

For the combined method, authentication edges receive a $1.5\times$ multiplier (\texttt{auth\_weight\_multiplier}), ensuring that the typically more informative authentication signals are not diluted by the larger number of flow edges. Self-loops and user-to-user edges are forced to a score of 0.0 and excluded from all downstream analysis.

\subsubsection{Path Enumeration and Scoring}
Beyond individual edges, we enumerate multi-hop paths through the graph to capture chain-like lateral movement patterns. Starting from each node, we perform breadth-first search up to a maximum of 4 hops, exploring only the top 10 outgoing edges per node ranked by edge score (pruning low-scored branches). Path enumeration is parallelized across 12 CPU workers using Python's \texttt{ProcessPoolExecutor}.

Each path is scored as the average of three aggregate measures of its constituent edge scores: geometric mean (emphasizing consistently anomalous paths), maximum edge score (capturing paths with at least one highly anomalous segment), and arithmetic mean. The top 50 paths by score are retained. After path enumeration, edge scores are boosted by adding 10\% of any path score that traverses the edge (\texttt{path\_boost\_factor} $= 0.1$), feeding path-level anomaly information back into edge-level detection.

\subsection{Threshold Optimization and Detection}

Anomalous edges are identified by applying a threshold to the final edge scores. Rather than using a fixed percentile, the pipeline uses an auto-optimization procedure: it sweeps across percentiles $[90, 95, 97, 99, 99.5, 99.9]$ and selects the threshold that maximizes the F1 score against the known red team pairs. This search is performed in pair-space: after thresholding, the set of anomalous source-destination pairs is compared against the ground-truth red team pairs to compute recall, precision, and F1. In our experiments, the auto-optimizer selected the 95th percentile for the auth-only method and the 97th percentile for both the flow-only and combined methods.

\subsection{Evaluation Metrics}

We evaluate each method using four metrics computed in pair-space: \emph{recall} (fraction of red team pairs detected), \emph{false positive rate} (fraction of non-red-team pairs flagged as anomalous), \emph{F1 score} (harmonic mean of precision and recall), and \emph{area under the ROC curve} (AUC). The AUC is computed using scikit-learn's \texttt{roc\_auc\_score} over all valid edges (excluding self-loops and user edges), with ground-truth labels derived from the red team pair set. This provides a threshold-independent measure of discriminative power.

\subsection{Baseline Methods}

To contextualize our graph-based approach, we compare against two classical unsupervised anomaly detection methods applied to the DAPT2020 dataset: One-Class SVM with an RBF kernel ($\gamma{=}\text{``scale''}$, $\nu{=}0.1$) and Isolation Forest (100 estimators, contamination$= 0.05$). Both are trained on normal-only samples and evaluated against held-out data containing both normal and lateral movement samples, operating on the 77 pre-extracted CICFlowMeter features. We note that this comparison has a methodological limitation: the baselines operate on DAPT2020 while our graph methods operate on LANL-2015, making direct comparison of absolute numbers invalid. Cross-dataset alignment is planned as future work.

\subsection{Implementation Details}

The pipeline is implemented in Python, using igraph for graph operations and scikit-learn for baseline methods and evaluation metrics. All hyperparameters are documented in a JSON configuration file (\texttt{pipeline\_config.json}), including scoring weights, path enumeration parameters, threshold search ranges, and baseline model settings. The streaming graph construction processes the full LANL-2015 dataset (135M+ events within the time windows) in a single pass, requiring approximately 4 GB of RAM. End-to-end pipeline execution for all three graph variants plus baselines completes in approximately 4 hours on a 12-core machine.